% 原文 pp.163–172（印刷 pp.1281–1290）
\section{Lecture 12：超对称场论}

\begin{keytakeaways}
本讲引入超对称（supersymmetry）场论：对称群非平凡地扩张为超群，其 Lie 超代数的奇元素给出费米型的无穷小对称。全讲的核心物理图像是：经典场方程的解若在某些（而非全部）超对称下不变，便自动在给定拓扑类中取能量的极小，这类解称为 BPS 态。核心数学工具是超空间（superspace）上的超场与超势（superpotential），以及把超对称条件化为一阶（Bogomolny 型）方程：在二维 $N=2$ 模型中，超对称孤子恰是超势 $W$ 的临界点之间沿 $h=\mathrm{Re}(e^{-i\alpha}W)$ 的梯度流线，其能量恰为 $|T|$，$T=W(\phi(\infty))-W(\phi(-\infty))$。四维 $N=1$、$N=2$ 的超庞加莱表示由相应 Clifford 代数二次型的秩决定（无质量表示维数 $2,4$，有质量表示维数 $4,16$），由此得到矢量多重态、手征多重态与超多重态等粒子内容。主要结论是：超对称代数可以含有中心荷（central charge），中心荷使 BPS 界 $M\ge|S|$ 能够被饱和；$N=2$ 规范理论中的中心荷必然是电荷与磁荷的线性组合，而质量恰等于 $|S|$ 的 BPS 磁单极子，是学期末研究 Seiberg--Witten 理论与 Donaldson 理论的出发点。
\end{keytakeaways}

% ===== 原文 p.163 =====
\subsection{超对称位形}
% 原文 §12.1：Supersymmetric Configurations

从今天起，我们研究具有超对称的场论，即其对称群可以非平凡地扩张为一个超群（supergroup）的场论。
按照惯例，所谓偶（even）与奇（odd）的（无穷小）对称，指的正是这个超群的 Lie 超代数（Lie superalgebra）中的偶元素与奇元素。

经常出现的情形是：经典场方程的一个解带有非平凡的奇对称。例如：

\begin{quote}
\noindent 梯度流线（gradient flow lines，Morse 理论中的）\\
\noindent 全纯曲线（holomorphic curves）\\
\noindent 规范理论瞬子（gauge theory instantons）\\
\noindent 磁单极子（monopoles）\\
\noindent Seiberg--Witten 解\\
\noindent 超 Kähler 结构（hyperKähler structures）\\
\noindent Calabi--Yau 度规\\
\noindent $G_2$ 与 $Spin_7$ 和乐的度规
\end{quote}

这些例子听起来差别很大，但它们有一个共同点：它们都是「在某些超对称下不变的经典场」的例子。
它们在数学上被大量研究，而在理解超对称量子系统的行为时也起着重要作用。

\subsection{超对称孤子（BPS 态）}
% 原文 §12.2：Supersymmetric solitons (BPS states)

首先考虑一个会出现 Morse 理论梯度流线的情形。
设 $h:\mathbb{R}^n\to\mathbb{R}$ 是 Morse 函数（即其临界点非退化，且 $|\nabla h|$ 在无穷远处增长）。Morse 函数总是只有有限多个临界点。
考虑（在闵可夫斯基号差下）映射 $\Phi:\mathbb{R}^{2|2}\to\mathbb{R}^n$ 的理论，其拉格朗日量为
\begin{equation}
  \tag{12.1}
  \label{eq:12-1}
  \mathcal{L}=\int d^2x\,d^2\theta\left(\frac12D_+\Phi D_-\Phi-h(\Phi)\right),
\end{equation}

\noindent\textit{（本讲笔记由 Pavel Etingof 与 David Kazhdan 整理。原文 p.163 页脚）}

% ===== 原文 p.164 =====
其中 $D_\pm=\dfrac{\partial}{\partial\theta_\pm}-\theta_\pm\partial_\pm$，$\partial_\pm=\dfrac{\partial}{\partial x^\pm}$，$x_\pm=\tfrac12(t\pm x)$。
此模型在 $Q_\pm=\dfrac{\partial}{\partial\theta_\pm}+\theta_\pm\partial_\pm$ 之下具有显然的超对称性。
这些超对称算符满足显然的对易关系 $Q_\pm^2=\partial_\pm$、$\{Q_+,Q_-\}=0$。

我们有 $(Q_+\pm Q_-)^2=2\dfrac{\partial}{\partial t}$。
因此，若我们要寻找具有时间平移对称性的经典解（即孤子），我们可以特别地在那些在某个超对称（比如说 $Q_++Q_-$ 或 $Q_+-Q_-$）下不变的解中去寻找。

我们有
\begin{equation}
  \tag{12.2}
  \label{eq:12-2}
  Q_+\pm Q_-=\left(\frac{\partial}{\partial\theta_+}\pm\frac{\partial}{\partial\theta_-}\right)
  +(\theta_+\pm\theta_-)\frac{\partial}{\partial t}+(\theta_+\mp\theta_-)\frac{\partial}{\partial x}.
\end{equation}

% [Tutor Note: (12.2) 的两组符号必须「配对反向」：$t$ 前的系数取 $(\theta_+\pm\theta_-)$，$x$ 前的系数就取 $(\theta_+\mp\theta_-)$。理由是由 $x_\pm=\tfrac12(t\pm x)$ 得 $\partial_+=\partial_t+\partial_x$、$\partial_-=\partial_t-\partial_x$，故 $Q_++Q_-$ 的 $x$ 系数为 $\theta_+-\theta_-$，$Q_+-Q_-$ 的 $x$ 系数为 $\theta_++\theta_-$——恰与 $t$ 系数反号。转写时若把两处都写成 $\pm$（同一符号），把 (12.3) 作用到 $\Phi=\phi+\theta_+\theta_-\nabla h(\phi)$ 上就会得到 $\frac{d\phi}{dx}=\nabla h$ 与 $\frac{d\phi}{dx}=-\nabla h$ 两个互相矛盾的方程，而不是同一个 (12.4)：这正是检验符号是否抄对的最快办法。]

因此，与时间无关的解的超对称条件是
\begin{equation}
  \tag{12.3}
  \label{eq:12-3}
  \left[\left(\frac{\partial}{\partial\theta_+}\pm\frac{\partial}{\partial\theta_-}\right)
  +(\theta_+\mp\theta_-)\frac{d}{dx}\right]\Phi=0.
\end{equation}

让我们寻找偶解。容易证明，这类解具有形式 $\Phi=\phi+\theta_+\theta_-\nabla h(\phi)$。对它们而言，超对称条件是
\begin{equation}
  \tag{12.4}
  \label{eq:12-4}
  \frac{d\phi}{dx}\mp\nabla h(\phi)=0,
\end{equation}
这正是梯度流线的条件。因此，超对称孤子就是梯度流的流线。

请注意，这些一阶方程蕴涵了二阶运动方程。事实上，容易证明运动方程为
\begin{equation}
  \tag{12.5}
  \label{eq:12-5}
  \partial_+\partial_-\phi+\nabla(\nabla h)^2(\phi)=0,
\end{equation}
或者对与时间无关的解，
\begin{equation}
  \tag{12.6}
  \label{eq:12-6}
  \frac{d^2\phi}{dx^2}=\nabla(\nabla h)^2(\phi),
\end{equation}
后者可由 (12.4) 对 $x$ 求导得到。

% [Tutor Note: 这里有一处**因子 2** 的疑点。把 (12.4) 即 $\phi'=\pm\nabla h$ 再对 $x$ 求导，得 $\phi''=\pm(\nabla h\cdot\nabla)\nabla h=\nabla\big(\tfrac12(\nabla h)^2\big)$，而不是 $\nabla\big((\nabla h)^2\big)=2(\nabla h\cdot\nabla)\nabla h$。也就是说，按 (12.7) 中势项 $\tfrac12(\nabla h)^2$ 的归一化，(12.5)–(12.6) 右端应带因子 $\tfrac12$；原文把 $\tfrac12$ 略去（或把记号 $\nabla(\nabla h)^2$ 另行归一化）。这一因子不影响「一阶方程蕴涵二阶方程」「超对称解饱和能量下界」这些定性结论，但若照字面把 (12.6) 与 (12.4) 联立验算会差 2 倍——可对照 (12.8) 中 $\tfrac12\big(\frac{d\phi}{dx}\mp\nabla h\big)^2$ 的配方来确认正确归一化。]

看清这一点的另一种方式如下：对与时间无关的场，其拉格朗日量（即哈密顿量）为
\begin{equation}
  \tag{12.7}
  \label{eq:12-7}
  H(\phi)=\int dx\left(\frac12\left(\frac{d\phi}{dx}\right)^2+\frac12(\nabla h(\phi))^2\right).
\end{equation}

改写 $H$，我们得到
\begin{equation}
  \tag{12.8}
  \label{eq:12-8}
  H(\phi)=\frac12\int dx\left(\frac{d\phi}{dx}\mp\nabla h(\phi)\right)^2
  \pm\int_{-\infty}^{\infty}dh(s).
\end{equation}

由于最后一项在有限能量的场空间上局部地为常数，一个超对称解在场空间的每个连通分支上都给出能量的全局极小。
该极小处的能量值为 $|S|$，其中 $S=\int dh$。

\textbf{定义.} 超对称孤子，即那些在某些但并非全部超对称下不变的经典解，称为经典 BPS 态（classical BPS states）。

% ===== 原文 p.165 =====
\subsection{BPS 态在量子理论中的作用}
% 原文 §12.3：The role of BPS states in quantum theory

我们在上面已经提到并使用了一个事实：向量场 $Q_+,Q_-$ 对易。
由于经典场方程的解空间是一个辛超流形（symplectic supermanifold），这些向量场必定（至少在局部上）由某些哈密顿函数 $\tilde Q_+,\tilde Q_-$ 生成（它们定义到相差一个局部常数函数）。
但这些函数不必 Poisson 对易：它们的 Poisson 括号必须是一个局部常数函数，而不一定等于零。

在我们的情形中，函数 $\tilde Q_+,\tilde Q_-$ 容易写出：若 $\Phi=\phi+\theta_+\psi_++\theta_-\psi_-+\theta_+\theta_-F$，则
\begin{equation}
  \tag{12.9}
  \label{eq:12-9}
  \tilde Q_+=\int dx\left(\psi_+\partial_+\phi+\psi_-\nabla h(\phi)\right),\qquad
  \tilde Q_-=\int dx\left(\psi_-\partial_-\phi-\psi_+\nabla h(\phi)\right).
\end{equation}

Poisson 括号的计算给出 $\{\tilde Q_+,\tilde Q_-\}=-2S$，其中 $S=\int dh=h(\phi(\infty))-h(\phi(-\infty))$。

% [Tutor Note: 由 $\{\tilde Q_+,\tilde Q_-\}=-2S$ 与 $\tilde Q_\pm^2=P_\pm$ 到 $(Q_+\pm Q_-)^2=2(H\mp\hat S)$ 之间，作者省掉了两步。第一步是纯代数：$(Q_+\pm Q_-)^2=Q_+^2+Q_-^2\pm\{Q_+,Q_-\}=P_++P_-\mp2S=2H\mp2S$，此处用到 $H=(P_++P_-)/2$。第二步是把 $S$ 提升为算符 $\hat S$：因为 $S$ 在场空间每个连通分支上是常数，$\hat S$ 属于超选择定则意义上的「中心」，可与一切局部算符对易，故在 $\mathcal H_a$ 上它取定标量值 $S_a$。再注意 $Q_+\pm Q_-$ 是厄米算符，其平方必非负，于是 $H\mp S_a\ge0$ 对两个符号都成立，这才给出作者下面那句「一般地 $H\ge|S_a|$」；原文只说「由于 $H\ge0$」，跳过了「平方非负」这一关键理由。]

从这个图像可以清楚看出 $Q_+,Q_-$ 这两个算符在量子理论中会变成什么。
我们将有 $Q_+^2=P_+$、$Q_-^2=P_-$（其中 $P_+,P_-$ 是相应的动量算符），并且 $(Q_+\pm Q_-)^2=2(H\mp\hat S)$，其中 $H=(P_++P_-)/2$ 是哈密顿量，而 $\hat S$（由 $S$ 量子化而得）与所有局部算符对易，因为 $S$ 在场空间的每个连通分支上是常数。

在量子理论中，场空间 $X$（即有限能量场构成的空间）的每个连通分支 $X_a$ 对应于 Hilbert 空间的一个直和项 $\mathcal{H}_a$。
在许多条件下，每一个超对称孤子 $\Phi\in X_a$ 都对应于 $\mathcal{H}_a$ 中一个同样是超对称的态 $\Psi$：$(Q_+\pm Q_-)\Psi=0$。（我们将讨论能够证明这样的 $\Psi$ 存在的例子。）
于是 $(H\mp S_a)\Psi=0$，其中 $S_a=\hat S|_{\mathcal{H}_a}$ 是一个标量。
因此，由于 $H\ge0$，我们有 $H\Psi=|S_a|\Psi$。
特别地，在这个只有两个超荷的超对称模型中，对 $X$ 的每个连通分支 $X_a$，至多存在一个超荷的线性组合（它由 $S_a$ 决定）能够湮灭 $\mathcal{H}_a$ 中的态。

一般地，在 $\mathcal{H}_a$ 上有 $H\ge|S_a|$。

现在我们要判断：是否存在一个超对称的量子态，即一个被 $Q_+\pm Q_-$ 湮灭、对应于该超对称经典态的量子态。
在微扰论的最低阶，人们发现除平移之外没有玻色零模，而恰有一个费米零模。
粗略地说，人们由此论证：一个费米零模不足以表示两个超对称，所以超对称的某个线性组合必定湮灭孤子的最低能量量子态。
这种被超对称的某个线性组合所湮灭的量子态，称为量子 BPS 态。

在这个情形下，要把「量子 BPS 态存在」的论证做得严格是笨拙的，因为孤子的费米零模数目为奇数。
$1+1$ 维中具有奇数个费米零模（且没有玻色零模）的孤子服从一种奇特的统计，既非玻色亦非费米；对此我们没有讨论过，我也不打算在此详述。\footnote{由 $n$ 个全同粒子组成的系统通常给出 $n$ 个对象的置换群 $S_n$ 的一个表示。对玻色子得到平凡表示；对费米子得到一维表示，它把置换 $\pi$ 映为 $(-1)^{\pi}$（即置换的符号）。具有奇数个费米零模的 $1+1$ 维孤子所表示的并不是 $S_n$，而是 $S_n$ 被 $\mathbb{Z}_2$ 中心扩张所得的群 $\hat S_n$。对偶数 $n$，其不可约表示的维数为 $2^{n/2}$。这些奇特的孤子还构成本学期公理化讲演中某些假设的反例，例如「量子场论的态在遥远的过去或未来渐近地约化为单粒子态的 Fock 空间」这一假设。不过我要向读者保证：这类反例显然只出现在 $1+1$ 维，以及（带有某些差别、涉及 Chern--Simons 相互作用与量子群的）$2+1$ 维。（原文脚注 1，跨 p.165--166 两页，此处合并排印。）}
只需指出：若有 $n$ 个彼此相距很远（相距

% ===== 原文 p.166 =====
远大于 $1/m$ 的距离，其中 $m$ 为质量间隙）的这类孤子，则对偶数 $n$，人们有 $n$ 个近似的费米零模（每个孤子附近支持一个这样的模），而 Clifford 代数的量子化给出 $2^{n/2}$ 个几乎简并的最低能态（在孤子位置固定的情形下）。
于是每个孤子的态数为 $\sqrt2$。
相反，若孤子不被超对称的任何线性组合所湮灭，那么在研究单孤子态时我们就必须表示超对称代数。
若 $H>|S|$，单个孤子的态将表示一个二维 Clifford 代数 $Q_+^2=Q_-^2=H$、$\{Q_+,Q_-\}=-2S$，它是非退化的，其最小不可约表示的维数为 2。
所以在此情形下，每个孤子将至少有 2 个态，而不只是 $\sqrt2$ 个。

证明量子 BPS 态存在的基本想法几乎总是一样的：数一数量子态的数目，并说明若这些态不被某些超对称所湮灭，数目就必须更大。
然而，目前的这个例子有点古怪，因为统计是奇特的，且每个粒子的量子态数目不是整数。
我不想彻底分析这个例子，因为所涉及的奇特之处在我们后续对量子场论的研究中不会再出现。
在随后的例子中，细节将更为直截了当。

\subsection{二维中的 $N=2$ 超对称}
% 原文 §12.4：N = 2 supersymmetry in 2 dimensions

考虑以 $x_+,x_-,\theta_+,\theta_-,\bar\theta_+,\bar\theta_-$ 为坐标的空间 $\mathbb{R}^{2|4}$。（关于此处所讲材料的更多细节，见 Supersolutions。）
此空间容许 $N=2$ 超对称代数的作用，超对称生成元为
\begin{equation}
  \tag{12.10}
  \label{eq:12-10}
  \begin{aligned}
    Q_+&=\frac{\partial}{\partial\theta_+}+\bar\theta_+\frac{\partial}{\partial x_+},
    &\bar Q_+&=\frac{\partial}{\partial\bar\theta_+}+\theta_+\frac{\partial}{\partial x_+},\\
    Q_-&=\frac{\partial}{\partial\theta_-}+\bar\theta_-\frac{\partial}{\partial x_-},
    &\bar Q_-&=\frac{\partial}{\partial\bar\theta_-}+\theta_-\frac{\partial}{\partial x_-},
  \end{aligned}
\end{equation}

我们还可以引入向量场
\begin{equation}
  \tag{12.11}
  \label{eq:12-11}
  \begin{aligned}
    D_+&=\frac{\partial}{\partial\theta_+}-\bar\theta_+\frac{\partial}{\partial x_+},
    &\bar D_+&=\frac{\partial}{\partial\bar\theta_+}-\theta_+\frac{\partial}{\partial x_+},\\
    D_-&=\frac{\partial}{\partial\theta_-}-\bar\theta_-\frac{\partial}{\partial x_-},
    &\bar D_-&=\frac{\partial}{\partial\bar\theta_-}-\theta_-\frac{\partial}{\partial x_-},
  \end{aligned}
\end{equation}
它们与超对称生成元对易。因此，任何用这些 $D$ 写出的拉格朗日量都是超对称的。

回忆一下：$\mathbb{R}^{2|4}$ 上的手征函数（chiral function，或手征超场）是满足方程 $\bar D_+\Phi=\bar D_-\Phi=0$ 的函数。
这些方程的一般解具有如下形式
\begin{equation}
  \tag{12.12}
  \label{eq:12-12}
  \begin{aligned}
    \Phi={}&\phi-\theta_+\bar\theta_+\partial_+\phi-\theta_-\bar\theta_-\partial_-\phi
    +\theta_+\bar\theta_+\theta_-\bar\theta_-\partial_+\partial_-\phi+\theta_+\theta_-F+{}\\
    &\theta_+\psi_++\theta_-\psi_--\theta_+\bar\theta_+\theta_-\partial_+\psi_-
    -\theta_-\bar\theta_-\theta_+\partial_-\psi_+.
  \end{aligned}
\end{equation}

% [Tutor Note: (12.12) 中四个费米项的符号与次序是纯粹由 Grassmann 代数决定的，极易抄错。要点有二：(1) 形如 $-\theta_+\bar\theta_+\theta_-\partial_+\psi_-$ 的项，其负号来自把 $\partial_+$ 作用在 $\bar\theta_+\theta_-\partial_+\phi$ 一类项上之后再移置 $\theta$ 变量；把 $\bar\theta_+$ 移过 $\theta_-$ 时必出 $(-1)$，漏掉这个负号会使整个展开不再是手征超场（可代回 $\bar D_\pm\Phi=0$ 检验）。(2) 末项印作 $\theta_-\bar\theta_-\theta_+\partial_-\psi_+$，而前面同类项印作 $\theta_+\bar\theta_+\theta_-\partial_+\psi_-$：作者的约定是「同一个 $\theta$ 与其共轭相邻、再乘另一个 $\theta$」，并不按 $\theta_+\theta_-\bar\theta_+\bar\theta_-$ 的字典序排列；翻译或重排时若擅自换序，必须同时补上相应的符号因子。]

% ===== 原文 p.167 =====
关于手征函数的更多内容，你可以在 Superhomework 与 Supersolutions 中读到。

考虑从 $\mathbb{R}^{2|4}$ 到 $\mathbb{C}^n$ 的手征映射 $\Phi$ 的理论，其拉格朗日量为
\begin{equation}
  \tag{12.13}
  \label{eq:12-13}
  \frac12\int d^2x\,d^4\theta\,\Phi\bar\Phi
  +\int d^2x\,d\theta_+d\theta_-\,W(\Phi)
  +\int d^2x\,d\bar\theta_+d\bar\theta_-\,\bar W(\Phi),
\end{equation}
其中 $W$ 是 $\mathbb{C}^n$ 上的全纯函数，称为超势（superpotential）。
此拉格朗日量是 $N=2$ 超对称的。
在分量形式下（取 $x_+=z$，$x_-=\bar z$），它形如
\begin{equation}
  \tag{12.14}
  \label{eq:12-14}
  \frac12\int d^2x\left(|d\phi|^2-|F|^2+W'(\phi)F+\bar W'(\phi)\bar F+\text{含费米子的项}\right).
\end{equation}

把全部费米子置零，并把「哑」场（dummy field）$F$ 取在驻点 $F=\bar W'(\phi)$，我们得到玻色能量泛函
\begin{equation}
  \tag{12.15}
  \label{eq:12-15}
  H=\frac12\int d^2x\left(|d\phi|^2+|W'(\phi)|^2\right).
\end{equation}

容易看出，对函数 $h=\mathrm{Re}(e^{-i\alpha}W)$（其中 $\alpha$ 为任意实数），此泛函与 (12.7) 相合。
因此，在经典层面我们所做的正是前一问题的特例。
然而，现在我们拥有更多的超对称，从而得到一个更有趣的理论。

% [Tutor Note: (12.13)、(12.14) 里 $\bar W$ 与 $\bar W'$ 的自变量都印成了 $\Phi$、$\phi$（而非 $\bar\Phi$、$\bar\phi$），(12.19)、(12.20) 亦然。这是作者对全纯量与反全纯量的宽松记号（$\bar W$ 视为 $\bar\phi$ 的函数，$\bar W'$ 视为对 $\bar\phi$ 求导）；计算时务必记住 $\bar W'(\phi)$ 实际是 $\overline{W'(\phi)}$，否则 (12.19) 与 (12.20) 之间「$\bar T=e^{-i\alpha}\int|W'|^2dx$」这一步就无法对上。]

\subsection{$N=2$ 的 BPS 态}
% 原文 §12.5：N = 2 BPS states

$N=2$ 超对称代数的对易关系为
\begin{equation}
  \tag{12.16}
  \label{eq:12-16}
  \begin{aligned}
    \{Q_\pm,\bar Q_\pm\}&=2P_\pm,\\
    \{Q_+,\bar Q_-\}&=\{\bar Q_+,Q_-\}=0\qquad(\text{由 }R\text{ 对称性}),\\
    \{Q_+,Q_-\}&=T,\qquad\{\bar Q_+,\bar Q_-\}=\bar T,
  \end{aligned}
\end{equation}
其中 $T$ 类比于 §12.2 中的 $S$——它是经典解空间上的局部常数函数（与 $S$ 不同，$T$ 不一定为实数）；为简洁起见，我们略去了 $Q$ 上的波浪号。
此外，在上述基下，所有 $Q$ 的平方都为零。

事实上，函数 $T$ 容易计算，就像前一问题中的函数 $S$ 一样。
与之前本质相同的计算给出
\begin{equation}
  \tag{12.17}
  \label{eq:12-17}
  T=W(\phi(\infty))-W(\phi(-\infty)).
\end{equation}

现在我们来考察超对称态。
取一个实数 $\alpha$，寻找在如下两个超对称下不变的态：$Q_1(\alpha)=Q_++e^{i\alpha}\bar Q_-$，$Q_2(\alpha)=\bar Q_+-e^{-i\alpha}Q_-$。
我们有
\begin{equation}
  \tag{12.18}
  \label{eq:12-18}
  \{Q_1,Q_2\}=2i\left(H-\mathrm{Re}(e^{-i\alpha}T)\right).
\end{equation}
这意味着超对称的经典态与时间无关。（由 $Q_1$、$Q_2$ 与 $T$ 生成的辛向量场，对此类态全都为零，故对 $H$ 也必定如此。）

% [Tutor Note: (12.18) 的前因子与符号需谨慎。若直接用 (12.16) 的代数按 $\{A,B\}$ 的双线性展开，得到的是 $\{Q_1,Q_2\}=2P_+-e^{-i\alpha}T+e^{i\alpha}\bar T-\bar T$，并不自动等于 $2i(H-\mathrm{Re}(e^{-i\alpha}T))$：其中既有 $P_+$ 与 $P_-$ 的不对称，又有一个显式的 $i$。合理的解释是作者在此对超荷采用了含 $i$ 的厄米归一化（并把 $P_-$ 一项吸收进 $H$ 的定义），或者 (12.18) 有印刷讹误。好在后续推导并不依赖 (12.18) 的精确前因子：由 $Q_1\Phi=Q_2\Phi=0$ 得 $\{Q_1,Q_2\}\Phi=0$，从而 $H=\mathrm{Re}(e^{-i\alpha}T)$，再取 $\alpha=\arg T$ 即得 (12.21) 的 $H=|T|$。阅读时把 (12.18) 当作「定义 $H$ 与 $T$ 之组合」的关系式即可。]

% ===== 原文 p.168 =====
对于偶（even，即玻色型）函数 $\Phi$，方程 $Q_1\Phi=0$ 给出（在与时间无关的情形下）：
\begin{equation}
  \tag{12.19}
  \label{eq:12-19}
  \frac{d\phi}{dx}=e^{-i\alpha}\,\overline{W'(\phi)},
\end{equation}
而第二个方程给出同样的结果。这特别蕴涵
\begin{equation}
  \tag{12.20}
  \label{eq:12-20}
  \bar T=\bar W(\phi(\infty))-\bar W(\phi(-\infty))
  =\int\bar W'(\phi)\frac{d\phi}{dx}dx
  =e^{-i\alpha}\int_{-\infty}^{\infty}|W'(\phi)|^2dx,
\end{equation}
从而 $\alpha=\arg T$。
于是由方程 (12.18) 我们得到：对于经典方程的超对称解，
\begin{equation}
  \tag{12.21}
  \label{eq:12-21}
  H=\mathrm{Re}(e^{-i\alpha}T)=|T|.
\end{equation}
对该解所在连通分支中的其他态，我们有 $H\ge|T|$。

% [Tutor Note: (12.20) 原文的积分下限印作 $\infty$（$\int_\infty^\infty$），显然应为 $-\infty$，本译文按 $-\infty$ 排印。此处的推导用到了两件事：一是 $\bar W(\phi)$ 沿 $x$ 的全导数 $\frac{d}{dx}\bar W=\bar W'\frac{d\phi}{dx}$，二是 (12.19) 把 $\frac{d\phi}{dx}$ 换成 $e^{-i\alpha}\overline{W'(\phi)}$，从而被积函数成为 $e^{-i\alpha}|W'|^2$——非负实函数乘上相因子，这正是「$H=|T|$ 且 $T$ 的相位可被 $\alpha$ 抵消」的根源，也是 BPS 界饱和的解析表现。]

那么，这些（有限能量的）超对称解在几何上究竟是什么？
由方程 (12.19) 清楚可见，它们是 $W$ 的临界点之间、沿 $h=\mathrm{Re}(e^{-i\alpha}W)$ 的梯度流的分界线（separatrices）。

现在转向量子理论。
由经典考虑我们看到，在我们的理论中 $H\ge|T|$，而在非退化情形下 $H=|T|$ 只对 BPS 态成立。
因此我们应当期望：在量子理论中发生同样的情形，只是 $T$ 先验地要取修正后的值。

考虑一般情形：经典 BPS 态附近哈密顿量的所有零模都来自超庞加莱群。
若对势的任意三个临界点 $a,b,c$ 都有 $|T_{ac}|<|T_{ab}|+|T_{bc}|$，则属于这种情形，其中 $T_{ab}$ 是 $T$ 在「从 $a$ 到 $b$ 的解空间分支」上的取值。
换言之，当 $a$ 与 $b$ 之间的梯度流线从不经过 $c$ 时即是如此。
因此，这对一般的 $W$ 成立，但在 $W$ 的族中，该条件会在实余维一的地方失效。

在量子理论中，我们预期当耦合趋于零时，$H$ 的最低本征值收敛到 $|T|$。
$H$ 最低本征值的重数（对给定的动量）由零模的量子化求出。
按我们的非退化假设，除平移之外没有玻色零模（如我们从第 9 讲所知，平移零模的存在与「孤子可以带有任意动量」这一事实相关），而只有两个费米零模，它们来自那些不保持该孤子不变的超对称。
量子化这两个费米零模，我们得到一个二维 Clifford 代数及其二维表示。
在弱耦合极限下，由此量子化得到的孤子谱是双曲面（hyperboloid）上半部分上某个等变向量丛的截面空间，其纤维是 Clifford 代数的标准二维不可约表示。

相反，若孤子的基态有 $H>|T|$，它就必须表示由全部四个 $Q$ 生成的四维 Clifford 代数，而基态的重数将是四。
换言之，孤子的态将是双曲面上一个秩为四的丛的截面。

所以，仅仅通过数出只有两个费米零模、因而量子孤子只有两个态，我们就能推断：孤子的基态必定是 BPS 的，且质量 $H=|T|$。
这当然是 §12.3 末尾所作论证的一个更直截了当的版本。
它不需要我们离开通常的玻色与费米统计的世界！

% ===== 原文 p.169 =====
\subsection{四维中的 $N=1$ 超对称}
% 原文 §12.6：N = 1 Supersymmetry in 4 dimensions

现在考虑四维中的超对称。（详见 Supersolutions 以及 Superhomework 的补编（Addendum）。）
我们从 $N=1$ 超对称开始。
在此情形下，超对称代数的奇部与二维中 $N=2$ 的情形类似。它由 $Q_\pm,\bar Q_\pm$ 生成，关系为
\begin{equation}
  \tag{12.22}
  \label{eq:12-22}
  \begin{aligned}
    \{Q_\alpha,\bar Q_{\dot\beta}\}&=2P_{\alpha\dot\beta},\\
    \{Q_\alpha,Q_\beta\}&=0,\\
    \{\bar Q_{\dot\alpha},\bar Q_{\dot\beta}\}&=0.
  \end{aligned}
\end{equation}
其中 $\alpha,\beta\in\{+,-\}$，而 $P_{\alpha\dot\beta}$ 是动量算符。（由于 Lorentz 表示的同构 $S_+\otimes S_-\to V_{\mathbb{C}}$，其中 $S_\pm$ 是旋量表示，四维中的动量算符在 $S_+\otimes S_-$ 中变换。）

% [Tutor Note: 点指标（dotted index）在这里不是装饰：$\bar Q$ 带点指标而 $Q$ 带无点指标，二者分属 $SL(2,\mathbb{C})$ 的两个不等价二分量旋量表示 $S_-$ 与 $S_+$。因此 $\{Q_\alpha,Q_\beta\}=0$ 与 $\{\bar Q_{\dot\alpha},\bar Q_{\dot\beta}\}=0$ 是两条独立的方程，而混合括号 $\{Q_\alpha,\bar Q_{\dot\beta}\}\propto P_{\alpha\dot\beta}$ 恰好把 $S_+\otimes S_-\cong V_{\mathbb{C}}$（四维矢量表示）显出来——这就是作者紧接着那段括注的内容。转写时把点漏掉（写成 $P_{\alpha\beta}$）会让这个张量结构消失，也会使 (12.22) 与 (12.24) 的指标配对无法核对。]

像 (12.16) 那样的中心扩张在这里不可能出现，因为它为庞加莱对称性所禁止（这种中心扩张必须处于 $SU(2)$ 的表示 $S^2\mathbb{C}^2$ 之中，其中 $\mathbb{C}^2$ 是标准表示；这是自旋 1 表示，不含任何不变量）。

% [Tutor Note: 这段群论依据值得展开。$\{Q_\alpha,Q_\beta\}$ 关于 $\alpha,\beta$ 必定对称（反对称部分当 $\alpha=\beta$ 时自动为零），故它落在 $S^2S_+\cong$ 自对偶 2-形式（即 $SU(2)$ 的自旋 1 表示 $S^2\mathbb{C}^2$）中；而自旋 1 表示不含 $SU(2)$ 不变量，所以不存在与 Lorentz 群对易的中心荷。$N=2$ 时情况改变：此时有两份超荷 $\alpha$ 指标可以配上反对称张量 $\epsilon_{\alpha\beta}$，$\{Q^{(1)}_\alpha,Q^{(2)}_\beta\}=\epsilon_{\alpha\beta}S$ 才成为可能（见 (12.24)）——换句话说，中心荷的存在依赖于「至少两份超对称」，这正是四维 $N=2$ 与 $N=1$ 的本质差别。]

对应于 Clifford 代数 (12.22) 的二次型，其行列式为 $(P^2)^2$。
当 $P^2=0$ 但 $P\ne0$ 时，此型的秩为 2；而当 $P^2\ne0$ 时，其秩为 4。
因此，超庞加莱群的无质量表示对应于秩 2 的 Clifford 代数的二维表示，而有质量表示对应于秩 4 的 Clifford 代数的四维表示。

更精确地说，考虑超庞加莱群的一个表示 $W$，以及其中 $P=p$ 的子空间 $W_p$，这里 $p\in\mathbb{R}^{1,3}$。
设 $G_p$ 是 $p$ 在旋转群中的稳定化子，$H_p$ 是 $G_p$ 中的极大环面（它总同构于 $U(1)$）。
关于 $H_p$，$W$ 分解为整数自旋与半整数自旋表示的直和。
这些自旋称为 $W$ 的螺旋度（helicity），每个螺旋度带有一个重数。

容易看出，超对称算符 $Q_\alpha,\bar Q_{\dot\alpha}$ 把螺旋度提高 $1/2$（因为它们处于 Lorentz 群的旋量表示之中）。
于是，无质量表示的螺旋度为 $j,j+1/2$，重数为 1；有质量表示的螺旋度为 $j,j+1/2,j+1$，重数为 $1,2,1$。

现在考虑 $N=1$ 超对称所允许的无质量粒子多重态。
这样的基本多重态有两个。

\begin{enumerate}
\item \textbf{矢量多重态（vector multiplet）}：一个规范场 $A$ 与一个处于伴随表示的手征旋量 $\lambda$。
  这是规范场的 $N=1$ 类比。
  特别地，在 $U(1)$ 情形下，这种场的最小理论是自由的。
  它含一个无质量矢量与一个无质量旋量，故螺旋度为 $-1,1$（对矢量）与 $-1/2,1/2$（对旋量）。
  特别地，此理论中出现的超庞加莱群的无质量表示是可约的，分裂为两个：螺旋度为 $-1,-1/2$ 的一个，与螺旋度为 $1/2,1$ 的一个。
  然而在 $\mathbb{R}$ 上这种分裂不存在，我们的表示是不可约的。
  这一陈述更物理的表述是：螺旋度 $-1,-1/2$ 的态与螺旋度 $1/2,1$ 的态通过 CPT 共轭相联系，所以一对若存在，另一对也必须存在。
  因此，这个表示（或物理学家所称的多重态）是含有规范玻色子的最小超对称表示。
\end{enumerate}

% ===== 原文 p.170 =====
\begin{enumerate}
\setcounter{enumi}{1}
\item \textbf{手征多重态（chiral multiplet）}：一个无质量复标量 $\phi$ 与一个无质量手征旋量 $\chi$。
  在此情形下情况是类似的。
  所得到的无质量表示含一个旋量与一个复标量，故螺旋度为 $0,0$（对标量）与 $-1/2,1/2$（对旋量）。
  因此这个表示又是两个表示之和，螺旋度分别为 $-1/2,0$ 与 $0,1/2$。
  此分解只在 $\mathbb{C}$ 上成立，而不在 $\mathbb{R}$ 上成立；在 $\mathbb{R}$ 上表示是不可约的，所以这个组合是含有标量的最小超对称（表示）。
\end{enumerate}

这些多重态的有质量版本如下。

\begin{enumerate}
\item \textbf{有质量矢量多重态。}
  含有有质量矢量的最小实超对称表示的螺旋度为 $-1,-1/2,0,1/2,1$，重数为 $1,2,2,2,1$。
  相应的 Clifford 代数的八维表示（在 $\mathbb{C}$ 上）是两个四维表示之和：仅超对称本身会给出螺旋度 $0,1/2,1$、重数 $1,2,1$ 的不可约多重态；CPT 给出其余的态。
  此多重态含一个有质量矢量、两个有质量旋量（手征与反手征）、一个实的有质量标量。
  可以把无质量矢量多重态与无质量手征多重态组合起来，并通过加入质量项（它可来自 Higgs 机制）作微扰而得到它。
\item \textbf{有质量手征多重态。}
  与无质量手征多重态相同（我们可以考虑同样的场但带上质量，使玻色子的质量等于费米子的质量；质量项可在超空间（superspace）中由超势里的二次项表示）。
\end{enumerate}

\subsection{四维中的 $N=2$ 超对称}
% 原文 §12.7：N = 2 Supersymmetry in 4 dimensions

在 $N=2$ 情形下，我们有两份算符 $Q$：$Q^{(1)}$ 与 $Q^{(2)}$。（同样，关于下文更多细节见 Supersolutions。）
$Q^{(1)}$ 与 $Q^{(1)}$ 或 $Q^{(2)}$ 与 $Q^{(2)}$ 的反交换子，生成我们刚才研究过的同一个 $N=1$ 超对称代数。
但在极小的 $N=2$ 代数中，$\{Q^{(1)},Q^{(2)}\}=0$。
然而现在有了中心扩张的可能性：它不再为庞加莱对称性所禁止。
中心扩张的特征是：$\{Q^{(1)},Q^{(2)}\}$ 非零（且为居中元）。

先考虑中心项为零的情形。
此时 Clifford 代数的二次型在有质量情形下（8 维中）非退化，在无质量情形下秩为 4。
因此，不可约的有质量表示应为 16 维，不可约的无质量表示应为 4 维。
表示的螺旋度像 $N=1$ 情形那样得到。
特别地，在有质量表示中，螺旋度按 $(j,j+1/2,j+1,j+3/2,j+2)$ 分组，重数为 $(1,4,6,4,1)$；在无质量表示中，螺旋度按 $(j,j+1/2,j+1)$ 分组，重数为 $(1,2,1)$。

现在考虑最简单的 $N=2$ 超对称理论。
它由（经典上）无质量的矢量多重态与伴随表示中的无质量手征多重态二者的场组合而成。
于是场为：来自矢量多重态的 $(A,\lambda)$，来自手征多重态的 $(\phi,\chi)$。
拉格朗日量是这些场上的极小 $N=2$ 超对称拉格朗日量。
这样的拉格朗日量存在，并由极小性条件唯一确定。
在 $U(1)$ 情形下理论是自由的，但在非阿贝尔情形下有非平凡的相互作用。

这个理论在红外做什么？在 $U(1)$ 情形下答案很简单：我们应把矢量多重态与手征多重态的表示相加。我们

% ===== 原文 p.171 =====
得到螺旋度 $(-1,-1/2,0,1/2,1)$，重数为 $(1,2,2,2,1)$。
这是 $N=2$ Clifford 代数的两个复共轭无质量表示之和：一个的螺旋度为 $(0,1/2,1)$、重数为 $(1,2,1)$，另一个的螺旋度为 $(-1,-1/2,0)$、重数为 $(1,2,1)$。

但现在让我们考虑非阿贝尔情形。
设规范群 $G$ 是单群。
玻色场是规范场 $A$ 与复化伴随表示 $\mathfrak{g}_{\mathbb{C}}$ 中的标量 $\phi$。
拉格朗日量的玻色部分为
\begin{equation}
  \tag{12.23}
  \label{eq:12-23}
  L_{\text{bosonic}}=\int d^4x\,\mathrm{Tr}\left(\frac{1}{4e^2}F^2+(d_A\phi)^2+[\phi,\bar\phi]^2\right).
\end{equation}

因此，经典基态空间的玻色部分是方程 $[\phi,\bar\phi]=0$ 的解集模去 $G$ 的作用。
此空间可等同于 $\mathfrak{t}_{\mathbb{C}}/W$，其中 $\mathfrak{t}$ 是 $\mathfrak{g}$ 中的极大交换子代数，$W$ 是 Weyl 群。
在 $G=SU(2)$ 的情形下，通过引入整体坐标 $u=\mathrm{Tr}(\phi^2)$（它参数化 Donaldson 理论的 $u$ 平面），此商被等同于 $\mathbb{C}$。

回忆第 2 讲：在此情形下我们有从 $G$ 到 $\phi$ 的中心化子 $H$ 的规范对称性破缺。
特别地，对 $u\ne0$ 的 $SU(2)$，规范对称性在经典上从 $SU(2)$ 破缺到 $U(1)$（Higgs 机制）。
关于残余 $U(1)$ 带非平凡电荷的规范场分量将获得质量。
由 $N=2$ 超对称，$\lambda,\phi,\psi$ 的相应分量也会如此。
于是，在 Hilbert 空间的电荷 $2$ 与 $-2$ 扇区中，最低能态将处于一个有质量表示之中，其螺旋度如前：$(-1,-1/2,0,1/2,1)$，重数为 $(1,2,2,2,1)$。
然而，没有中心扩张就不存在这样的表示。
因此，没有中心扩张我们便得到矛盾，从而中心扩张必然出现。

中心扩张将表现在 $Q^{(1)}_\alpha$ 与 $Q^{(2)}_\beta$ 的对易关系之中：
\begin{equation}
  \tag{12.24}
  \label{eq:12-24}
  \{Q^{(1)}_\alpha,Q^{(2)}_\beta\}=\epsilon_{\alpha\beta}S,
\end{equation}
其中 $S$ 是与所有局部算符对易的算符（它称为中心荷（central charge））。
在没有中心荷时，代数具有 $U(2)$ 的 R 对称性（作用于指标 1 与 2），但中心荷将此对称性破缺为 $SU(2)$。
无论如何，在量子层面，我们理论中的手征 $U(1)$ 对称性是反常的，因为费米子测度不是不变的；反常把手征 $U(1)$ 对称性破缺为 $\mathbb{Z}/8\mathbb{Z}$。

% [Tutor Note: 末句的 $\mathbb{Z}/8\mathbb{Z}$ 值得核对再引用。按 $N=1$ 语言，含一个伴随表示费米子的 $SU(N)$ 理论中，连续 R 对称性通常只破缺到 $\mathbb{Z}_{2N}$，对 $SU(2)$ 即 $\mathbb{Z}_4$；而这里讨论的是 $N=2$ 理论的 $U(1)_R$（$U(2)_R=SU(2)_R\times U(1)_R$ 中的那个 $U(1)$），其生成元对两个超荷的荷分配与 $N=1$ 归一化不同，因此残留的离散群指数会随归一化改变。作者在此写 $\mathbb{Z}_8$，读者引用时务必注明「按本讲的 $U(1)_R$ 归一化」，不要直接与 Seiberg 原文的 $\mathbb{Z}_{2N_c}$ 混用；原文此处还漏了一个词（"breaks the chiral $U(1)$ symmetry is to $Z/8Z$"），本译文按语义补齐为「破缺为」。]

在此情形下，上面考虑的有质量粒子，其质量必须恰为 $|S|$。
这正是使 Clifford 代数的二次型退化的那个值；而正如 §12.5 讨论的情形，若质量的取值不同，粒子数就必须更大。

此理论还有 BPS 孤子——磁单极子。
若我们处于真空 $u\in\mathbb{C}$，则有
\begin{equation}
  \tag{12.25}
  \label{eq:12-25}
  \phi\sim\frac{1}{\sqrt2}\begin{pmatrix}a&0\\0&-a\end{pmatrix}.
\end{equation}
其中 $a=\pm u^{1/2}$。

% [Tutor Note: (12.25) 中 $\frac{1}{\sqrt2}$ 这个看似随意的归一化，正是为了让 $u=\mathrm{Tr}(\phi^2)$ 与 $a$ 的关系取最简形式：把 $\phi=\frac{1}{\sqrt2}\mathrm{diag}(a,-a)$ 代入，得 $u=\mathrm{Tr}(\phi^2)=\frac12\big(a^2+(-a)^2\big)=a^2$，故 $a=\pm u^{1/2}$——这正好解释了紧接着的那句「其中 $a=\pm u^{1/2}$」。若改用 $\phi=\mathrm{diag}(a,-a)$ 的约定，则 $u=2a^2$、$a=\pm\sqrt{u/2}$，此后所有关于 $a$、中心荷 $S$ 与 BPS 质量 $M=|S|$ 的表达式都会整体差一个 $\sqrt2$。比对不同文献的 Seiberg--Witten 公式（$u=\langle\mathrm{Tr}\,\phi^2\rangle$、$a_D=\partial\mathcal{F}/\partial a$）时，这是最常见的因子来源。]

设 $a=\rho e^{i\alpha}$，$\rho$ 为实数。
我们应寻找满足无穷远处条件 (12.25) 的经典 BPS 解（即与时间无关、且在某些——事实上是一半——超对称下不变的解）。若

% ===== 原文 p.172 =====
$S=|S|e^{i\beta}$，则「在一半超对称下不变」这一方程是
\begin{equation}
  \tag{12.26}
  \label{eq:12-26}
  F=e^{-i\beta}*d_A\phi,
\end{equation}
其中 $*$ 是 $\mathbb{R}^3$ 中的 Hodge 星。（这就是 BPS 磁单极子方程。）
由于 $F$ 为实，我们必须有 $\alpha=\beta$ 模 $\pi$。

可以证明这样的 BPS 解存在；事实上关于它们有一套广泛的数学理论。
它们由磁荷的值（或无穷远处的第一 Chern 类）刻画，后者已在第 9 讲中引入。
对磁荷 1，解空间如下。
约化空间同构于 $\mathbb{R}^3\times S^1$，理由已在第 9 讲中解释。（回忆 $\mathbb{R}^3\times S^1$ 曾作为「由空间平移与无穷远处的 $U(1)$ 电荷旋转所构成的群」的齐性空间出现。）
该解在四个超对称的线性组合下不变；恰有四个费米零模，它们由其余的超对称生成。
量子化这些费米零模，给出一个四维 Clifford 代数及其四维不可约表示；因此电荷为 1 的磁单极子态处于超对称的四维表示之中。
此表示事实上是一个有质量超多重态（hypermultiplet）——两个自旋零的态与两个自旋二分之一的态。（反磁单极子由 CPT 与磁单极子相联系，同样有四个态，自旋相同。）
这些态必然是 BPS 的，质量 $M=|S|$；因为若 $M>|S|$，就必须表示一个八维 Clifford 代数，态数将是 16 而不是 4。
这些 BPS 态在我们于学期末分析此理论的动力学（及其在 Donaldson 理论中的应用）时，将起非常重要的作用。

\textbf{注（Remark）.} 用流（currents）计算 $Q$ 的对易子，可以证明经典地有
\begin{equation}
  \tag{12.27}
  \label{eq:12-27}
  S=\int_\Sigma\left(\phi*F+\frac{i}{e^2}\phi F\right),
\end{equation}
其中 $\Sigma$ 是空间切片中一个遥远的球面。
所以它是有效 $U(1)$ 理论中电荷与磁荷的线性组合。
事实上，$S$ 必然是电荷与磁荷的线性组合，原因很简单：此理论没有别的与庞加莱群对易的守恒荷。

% [Tutor Note: (12.27) 中的 $*$ 与 (12.26) 中的 $*$ 不是同一件事，极易混用。(12.26) 里的 $*$ 是**空间** $\mathbb{R}^3$ 内的 Hodge 星，把 1-形式 $d_A\phi$ 变成 2-形式 $*d_A\phi$，故 $F=e^{-i\beta}*d_A\phi$ 才可能是 2-形式与 2-形式的等式。(12.27) 里 $\Sigma$ 是空间切片中的 2-球面，$F$ 与 $*F$（四维 Hodge 星限制到空间 2-形式上，即电、磁对偶场强）都是 2-形式：$\int_\Sigma F$ 给出磁荷，$\int_\Sigma *F$ 给出电通量。此外 $\phi$ 是 Lie 代数值的，积分时隐含 $\mathrm{Tr}$ 与在 $\Sigma$ 上取 $\phi\to\phi_\infty$ 的渐近值；系数 $i/e^2$ 中的 $i$ 保证电、磁两部分以 $1/e^2$ 与 $\theta/2\pi$ 的方式组合成复中心荷。若把两处的 $*$ 混为一谈，$\phi*F$ 就会被误当作 1-形式而无法在 2-球面上积分。]

\subsection*{参考文献与交叉引用}
% 原文 References and cross references

§12.2 处理的 $\mathbb{R}^{2|2}$ 上的模型，在 [I-Supersolutions, §4.5] 中讨论。
具有 4 个超荷的超对称的基本公式，出现在 [I-Supersolutions, §§2.4, 2.6]。
§12.5 讨论的 $\mathbb{R}^{2|4}$ 模型，是 $\mathbb{R}^{4|4}$ 上模型的维数约化。
后者在 [I-Supersolutions] 第 5 章中讨论，不过中心扩张只在维数约化之后的模型中才看得到。
四维中各种超对称多重态，是 [I-Homework] 中 Superhomework 补编的主题。
四维中 $N=2$ 规范理论的量子行为，在本系列第 17、18 讲中讨论。

关于含孤子理论的超对称代数中中心荷的出现，见
\begin{quote}
\noindent E. Witten, D. Olive, \emph{Supersymmetry algebras that include topological charges}, Phys.\ Lett.\ 78B (1978), 97.
\end{quote}

关于四维 $N=2$ 超对称 Yang--Mills 理论中 BPS 磁单极子的数学理论，见
\begin{quote}
\noindent M. F. Atiyah, N. Hitchin, \emph{The Geometry And Dynamics Of Magnetic Monopoles}, M. B. Porter Lectures, 1988.
\end{quote}

% [Tutor Note: 原文的 [I-Supersolutions]、[I-Homework] 等是本丛书（\emph{Quantum Fields and Strings: A Course for Mathematicians}）第 I 卷内对应章节的文献代号；「Supersolutions」「Superhomework」指该卷第 5 章附近的超对称习题与解答材料，本讲多处（§12.4 的 $\mathbb{R}^{2|4}$、§12.6 的四维多重态、§12.7 的 $N=2$ 理论）都以它们为详细推导的补充，读者若不熟悉该卷，可对照本系列第 17、18 讲中 $N=2$ 规范理论动力学的完整处理。]
